\section*{章末练习}

\begin{longtable}{L{0.08\linewidth}L{0.32\linewidth}L{0.5\linewidth}}
\toprule
题号 & 类型 & 题目 \\
\midrule
1 & 单元读取 & 解释 DRAM 读取为何具有破坏性，以及感应放大器为何还承担恢复写回。 \\
2 & 时序 & 一次访问命中已打开 row 与访问另一 row 的命令序列有何差异？指出 ACT、READ/WRITE 与 PRE 的作用。 \\
3 & 地址映射 & 地址在 channel、rank、bank、row、column 间的映射如何影响并行性和 row-buffer 命中？ \\
4 & 刷新 & 刷新为何会形成不可用窗口？控制器怎样在正确性与尾延迟之间调度刷新？ \\
5 & DDR 训练 & 写 leveling、读采样训练和 per-bit deskew 分别在校准什么？为何不能只按原理图长度固定一次？ \\
6 & ECC & SECDED 可纠正错误与不可纠正错误到达 CPU 时应走哪些不同路径？ \\
\bottomrule
\end{longtable}

\section*{参考解答与要点}

\begin{longtable}{L{0.08\linewidth}L{0.32\linewidth}L{0.5\linewidth}}
\toprule
题号 & 类型 & 解答要点 \\
\midrule
1 & 单元读取 & 微小单元电荷与位线共享后改变原状态，感放把微小差分放大成全摆幅，并把判定值重新写回单元。 \\
2 & 时序 & row hit 可直接列访问；不同 row 需先 PRE 关闭旧 row，再 ACT 打开新 row，等待规定时序后 READ/WRITE。 \\
3 & 地址映射 & 把连续请求分散到不同 channel/bank 可增加并行，却可能减少同 row 局部性；集中 column 位有利于 row hit，但热点会集中。 \\
4 & 刷新 & 刷新期间相关 bank/rank 不能服务普通访问；控制器可提前、延后或细分刷新，但必须在保持时间约束内完成，并控制最长阻塞。 \\
5 & DDR 训练 & 分别校准命令/选通信号关系、接收采样窗和各数据 bit 偏斜；工艺、电压、温度及板级差异使静态名义长度不足以保证眼图中心。 \\
6 & ECC & 可纠正错误修正数据、记录并可 scrub；不可纠正错误不得作为普通数据交付，应报告机器检查、定位地址并隔离页面或故障部件。 \\
\bottomrule
\end{longtable}
